% Definierte Konstruktionsnotation auf Grundlage der Blattwortabbildung.
% Einbindung nach TreeLeafWordEquations und vor den Klammerungsmengen.

\subsection{Ein Baum klammert ein Wort}

Die Schreibweise \(\TreeBrackets{A}{T}{w}\) wird als
\glqq\(T\) klammert \(w\)\grqq\ gelesen. Sie bezeichnet die Zugehörigkeit
zum bereits konstruierten Funktionsgraphen des Blattwortes. Die folgenden
Sätze stellen dessen Typisierung und Konstruktionsgesetze für kurze
formale Ableitungen bereit.

\FormulaDefDeltaK[Ein Baum klammert ein Wort]{%
  \TreeBrackets{A}{T}{w}\coloneqq
    (T,w)\in\operatorname{wort}_A%
}{TreeBracketingRelationDef}{%
  \DeltaRow{Mengen}{A\dsep T\dsep w}
  \DeltaRow{Funktionen}{\operatorname{wort}_A}
  \DeltaRow{\textbf{Neues Symbol}}{}
  \DeltaPrem{Klammerungsrelation}{\TreeBrackets{A}{T}{w}}
}

\FormulaThmDeltaK[Baumtypisierung der Klammerungsrelation]{%
  \TreeBrackets{A}{T}{w}\vdash T\in\BracketTreeSet{A}%
}{TreeBracketingRelationTreeTyping}{%
  \DeltaRow{Mengen}{A\dsep T\dsep w}
}
\begin{tabproofwide}
  \proofstepwidestar[1]{\TreeBrackets{A}{T}{w}}{\rA}
  \proofstepwidestar[1]{(T,w)\in\operatorname{wort}_A}{%
    \FormulaRefAuto{TreeBracketingRelationDef}[def]{1}}
  \proofstepwidestar[]{%
    \operatorname{wort}_A\colon
      \BracketTreeSet{A}\to\NonemptyWords{A}}{%
    \FormulaRefAuto{TreeLeafWordTyping}[thm]}
  \proofstepwidestar[1]{T\in\BracketTreeSet{A}}{%
    \FormulaRefAuto{F\colon A\to B\dsep (x,y)\in F\vdash
      x\in A}[thm]{3,2}}
\end{tabproofwide}

\FormulaThmDeltaK[Worttypisierung der Klammerungsrelation]{%
  \TreeBrackets{A}{T}{w}\vdash w\in\NonemptyWords{A}%
}{TreeBracketingRelationWordTyping}{%
  \DeltaRow{Mengen}{A\dsep T\dsep w}
}
\begin{tabproofwide}
  \proofstepwidestar[1]{\TreeBrackets{A}{T}{w}}{\rA}
  \proofstepwidestar[1]{(T,w)\in\operatorname{wort}_A}{%
    \FormulaRefAuto{TreeBracketingRelationDef}[def]{1}}
  \proofstepwidestar[]{%
    \operatorname{wort}_A\colon
      \BracketTreeSet{A}\to\NonemptyWords{A}}{%
    \FormulaRefAuto{TreeLeafWordTyping}[thm]}
  \proofstepwidestar[1]{w\in\NonemptyWords{A}}{%
    \FormulaRefAuto{F\colon A\to B\dsep (x,y)\in F\vdash
      y\in B}[thm]{3,2}}
\end{tabproofwide}

\FormulaThmDeltaK[Auswertung der Klammerungsrelation]{%
  \TreeBrackets{A}{T}{w}\vdash\TreeLeafWord{A}{T}=w%
}{TreeBracketingRelationElimination}{%
  \DeltaRow{Mengen}{A\dsep T\dsep w}
  \DeltaRow{Funktionen}{\operatorname{wort}_A}
}
\begin{tabproofwide}
  \proofstepwidestar[1]{\TreeBrackets{A}{T}{w}}{\rA}
  \proofstepwidestar[1]{(T,w)\in\operatorname{wort}_A}{%
    \FormulaRefAuto{TreeBracketingRelationDef}[def]{1}}
  \proofstepwidestar[]{%
    \operatorname{wort}_A\colon
      \BracketTreeSet{A}\to\NonemptyWords{A}}{%
    \FormulaRefAuto{TreeLeafWordTyping}[thm]}
  \proofstepwidestar[1]{\TreeLeafWord{A}{T}=w}{%
    \FormulaRefAuto{F\colon A\to B\dsep (x,y)\in F\vdash
      F(x)=y}[thm]{3,2}}
\end{tabproofwide}

\FormulaThmDeltaK[Einführung der Klammerungsrelation]{%
  T\in\BracketTreeSet{A}\dsep\TreeLeafWord{A}{T}=w\vdash
    \TreeBrackets{A}{T}{w}%
}{TreeBracketingRelationIntroduction}{%
  \DeltaRow{Mengen}{A\dsep T\dsep w}
  \DeltaRow{Funktionen}{\operatorname{wort}_A}
}
\begin{tabproofwide}
  \proofstepwidestar[1]{T\in\BracketTreeSet{A}}{\rA}
  \proofstepwidestar[2]{\TreeLeafWord{A}{T}=w}{\rA}
  \proofstepwidestar[]{%
    \operatorname{wort}_A\colon
      \BracketTreeSet{A}\to\NonemptyWords{A}}{%
    \FormulaRefAuto{TreeLeafWordTyping}[thm]}
  \proofstepwidestar[1,2]{(T,w)\in\operatorname{wort}_A}{%
    \FormulaRefAuto{F\colon A\to B\dsep x\in A\dsep
      F(x)=y\vdash (x,y)\in F}[thm]{3,1,2}}
  \proofstepwidestar[1,2]{\TreeBrackets{A}{T}{w}}{%
    \FormulaRefAuto{TreeBracketingRelationDef}[def]{4}}
\end{tabproofwide}

\subsection{Blattregel und Knotenregel}

Ein Blatt klammert das zugehörige Buchstabenwort. Ein Knoten klammert
die Konkatenation der Wörter seiner beiden Teilbäume. Die Regeln
enthalten damit zugleich die benötigten Aussagen über den Baum und
sein Blattwort.

\FormulaThmDeltaK[Blattregel der Klammerungsrelation]{%
  a\in A\vdash
    \TreeBrackets{A}{\LeafTree{A}{a}}{\LetterWord{a}}%
}{TreeBracketingLeafIntroduction}{%
  \DeltaRow{Mengen}{A\dsep a}
}
\begin{tabproofwide}
  \proofstepwidestar[1]{a\in A}{\rA}
  \proofstepwidestar[1]{\LeafTree{A}{a}\in\BracketTreeSet{A}}{%
    \FormulaRefAuto{TreeLeafClosure}[thm]{1}}
  \proofstepwidestar[1]{%
    \TreeLeafWord{A}{\LeafTree{A}{a}}=\LetterWord{a}}{%
    \FormulaRefAuto{TreeLeafWordLeafEquation}[thm]{1}}
  \proofstepwidestar[1]{%
    \TreeBrackets{A}{\LeafTree{A}{a}}{\LetterWord{a}}}{%
    \FormulaRefAuto{TreeBracketingRelationIntroduction}[thm]{2,3}}
\end{tabproofwide}

\FormulaThmDeltaK[Knotenregel der Klammerungsrelation]{%
  \begin{aligned}[t]
    &\TreeBrackets{A}{S}{u}\dsep\TreeBrackets{A}{T}{v}\vdash{}\\[-2pt]
    &\TreeBrackets{A}{\NodeTree{A}{S}{T}}{\WordConcat{u}{v}}
  \end{aligned}%
}{TreeBracketingNodeIntroduction}{%
  \DeltaRow{Mengen}{A\dsep S\dsep T\dsep u\dsep v}
}
\begin{tabproofwide}
  \proofstepwidestar[1]{\TreeBrackets{A}{S}{u}}{\rA}
  \proofstepwidestar[2]{\TreeBrackets{A}{T}{v}}{\rA}
  \proofstepwidestar[1]{S\in\BracketTreeSet{A}}{%
    \FormulaRefAuto{TreeBracketingRelationTreeTyping}[thm]{1}}
  \proofstepwidestar[2]{T\in\BracketTreeSet{A}}{%
    \FormulaRefAuto{TreeBracketingRelationTreeTyping}[thm]{2}}
  \proofstepwidestar[1]{u\in\NonemptyWords{A}}{%
    \FormulaRefAuto{TreeBracketingRelationWordTyping}[thm]{1}}
  \proofstepwidestar[2]{v\in\NonemptyWords{A}}{%
    \FormulaRefAuto{TreeBracketingRelationWordTyping}[thm]{2}}
  \proofstepwidestar[1]{u\in\FiniteWords{A}}{%
    \FormulaRefAuto{NonemptyWordFinite}[thm]{5}}
  \proofstepwidestar[2]{v\in\FiniteWords{A}}{%
    \FormulaRefAuto{NonemptyWordFinite}[thm]{6}}
  \proofstepwidestar[1]{\TreeLeafWord{A}{S}=u}{%
    \FormulaRefAuto{TreeBracketingRelationElimination}[thm]{1}}
  \proofstepwidestar[2]{\TreeLeafWord{A}{T}=v}{%
    \FormulaRefAuto{TreeBracketingRelationElimination}[thm]{2}}
  \proofstepwidestar[1,2]{S,T\in\BracketTreeSet{A}}{\rAI{3,4}}
  \proofstepwidestar[1,2]{%
    \NodeTree{A}{S}{T}\in\BracketTreeSet{A}}{%
    \FormulaRefAuto{TreeNodeClosure}[thm]{11}}
  \proofstepwidestar[1,2]{%
    \begin{aligned}[t]
      &\TreeLeafWord{A}{\NodeTree{A}{S}{T}}\\[-2pt]
      &\quad=\WordConcat{\TreeLeafWord{A}{S}}{\TreeLeafWord{A}{T}}
    \end{aligned}}{%
    \FormulaRefAuto{TreeLeafWordNodeEquation}[thm]{11}}
  \proofstepwidestar[1,2]{%
    \begin{aligned}[t]
      &\TreeLeafWord{A}{\NodeTree{A}{S}{T}}\\[-2pt]
      &\quad=\WordConcat{u}{\TreeLeafWord{A}{T}}
    \end{aligned}}{\rIE{9,13}}
  \proofstepwidestar[1,2]{%
    \TreeLeafWord{A}{\NodeTree{A}{S}{T}}=\WordConcat{u}{v}}{%
    \rIE{10,14}}
  \proofstepwidestar[1,2]{%
    \TreeBrackets{A}{\NodeTree{A}{S}{T}}{\WordConcat{u}{v}}}{%
    \FormulaRefAuto{TreeBracketingRelationIntroduction}[thm]{12,15}}
\end{tabproofwide}

\subsection{Ein Blatt rechts anfügen}

Der folgende Operator bildet einen neuen Knoten mit dem bisherigen Baum
als linkem Teilbaum und einem neuen Blatt als rechtem Teilbaum. In der
Anfügeregel werden Baum und Wort dadurch in derselben Richtung erweitert.

\FormulaDefDeltaK[Rechtsanfügen eines Blattes]{%
  \begin{aligned}[t]
    &T\in\BracketTreeSet{A}\dsep a\in A\vdash{}\\[-2pt]
    &\TreeAppendLeaf{A}{T}{a}\coloneqq
      \NodeTree{A}{T}{\LeafTree{A}{a}}
  \end{aligned}%
}{TreeAppendLeafDef}{%
  \DeltaRow{Mengen}{A\dsep T\dsep a}
  \DeltaRow{\textbf{Neues Symbol}}{}
  \DeltaPrem{Blattanfügung}{\TreeAppendLeaf{A}{T}{a}}
}

\FormulaThmDeltaK[Die Blattanfügung ist ein Klammerungsbaum]{%
  T\in\BracketTreeSet{A}\dsep a\in A\vdash
    \TreeAppendLeaf{A}{T}{a}\in\BracketTreeSet{A}%
}{TreeAppendLeafTyping}{%
  \DeltaRow{Mengen}{A\dsep T\dsep a}
}
\begin{tabproofwide}
  \proofstepwidestar[1]{T\in\BracketTreeSet{A}}{\rA}
  \proofstepwidestar[2]{a\in A}{\rA}
  \proofstepwidestar[2]{\LeafTree{A}{a}\in\BracketTreeSet{A}}{%
    \FormulaRefAuto{TreeLeafClosure}[thm]{2}}
  \proofstepwidestar[1,2]{T,\LeafTree{A}{a}\in\BracketTreeSet{A}}{%
    \rAI{1,3}}
  \proofstepwidestar[1,2]{%
    \NodeTree{A}{T}{\LeafTree{A}{a}}\in\BracketTreeSet{A}}{%
    \FormulaRefAuto{TreeNodeClosure}[thm]{4}}
  \proofstepwidestar[1,2]{%
    \TreeAppendLeaf{A}{T}{a}=\NodeTree{A}{T}{\LeafTree{A}{a}}}{%
    \FormulaRefAuto{TreeAppendLeafDef}[def]{1,2}}
  \proofstepwidestar[1,2]{%
    \NodeTree{A}{T}{\LeafTree{A}{a}}=\TreeAppendLeaf{A}{T}{a}}{%
    \FormulaRefAuto{a=b\vdash b=a}[thm]{6}}
  \proofstepwidestar[1,2]{%
    \TreeAppendLeaf{A}{T}{a}\in\BracketTreeSet{A}}{\rIE{7,5}}
\end{tabproofwide}

\FormulaThmDeltaK[Anfügeregel der Klammerungsrelation]{%
  \begin{aligned}[t]
    &\TreeBrackets{A}{T}{u}\dsep a\in A\vdash{}\\[-2pt]
    &\TreeBrackets{A}{\TreeAppendLeaf{A}{T}{a}}
      {\WordConcat{u}{\LetterWord{a}}}
  \end{aligned}%
}{TreeBracketingAppendLeafIntroduction}{%
  \DeltaRow{Mengen}{A\dsep T\dsep u\dsep a}
}
\begin{tabproofwide}
  \proofstepwidestar[1]{\TreeBrackets{A}{T}{u}}{\rA}
  \proofstepwidestar[2]{a\in A}{\rA}
  \proofstepwidestar[1]{T\in\BracketTreeSet{A}}{%
    \FormulaRefAuto{TreeBracketingRelationTreeTyping}[thm]{1}}
  \proofstepwidestar[2]{%
    \TreeBrackets{A}{\LeafTree{A}{a}}{\LetterWord{a}}}{%
    \FormulaRefAuto{TreeBracketingLeafIntroduction}[thm]{2}}
  \proofstepwidestar[1,2]{%
    \TreeBrackets{A}{\NodeTree{A}{T}{\LeafTree{A}{a}}}
      {\WordConcat{u}{\LetterWord{a}}}}{%
    \FormulaRefAuto{TreeBracketingNodeIntroduction}[thm]{1,4}}
  \proofstepwidestar[1,2]{%
    \TreeAppendLeaf{A}{T}{a}=\NodeTree{A}{T}{\LeafTree{A}{a}}}{%
    \FormulaRefAuto{TreeAppendLeafDef}[def]{3,2}}
  \proofstepwidestar[1,2]{%
    \NodeTree{A}{T}{\LeafTree{A}{a}}=\TreeAppendLeaf{A}{T}{a}}{%
    \FormulaRefAuto{a=b\vdash b=a}[thm]{6}}
  \proofstepwidestar[1,2]{%
    \TreeBrackets{A}{\TreeAppendLeaf{A}{T}{a}}
      {\WordConcat{u}{\LetterWord{a}}}}{\rIE{7,5}}
\end{tabproofwide}
